% Part: first-order logic

\documentclass[../../include/open-logic-part]{subfiles}

\begin{document}

\olpart{fol}{પ્રથમ-ક્રમ તર્કશાસ્ત્ર}

\begin{editorial}
  આ ભાગમાં પ્રથમ-ક્રમ તર્કશાસ્ત્રના અધિસિદ્ધાંતને પૂર્ણતા સુધી આવરી
  લેવાયો છે. હાલ તે વિધાનાત્મક તર્કશાસ્ત્રની અલગ રજૂઆત પર આધાર રાખતો
  નથી; બધું અહીં સાબિત કરવામાં આવે છે. જો ``FOL'' ટૅગ false હોય, તો
  સ્રોત ફાઇલો પરિમાણકો વિશેની સામગ્રી કાઢી નાખશે (અને
  ``!!{structure}''ને ``!!{valuation}'', $\Struct{M}$ને $\pAssign{v}$
  વગેરેથી બદલશે). હકીકતમાં, વિધાનાત્મક તર્કશાસ્ત્રના ભાગની મોટાભાગની
  સામગ્રી એ ``FOL'' ટૅગ બંધ રાખેલું પ્રથમ-ક્રમનું જ રૂપ છે.

  વિધાનાત્મક તર્કશાસ્ત્રનો ભાગ સામેલ હોય, તો ઘણી સામગ્રીનું પુનરાવર્તન
  થાય છે. જોકે આ ભાગ વિધાનાત્મક તર્કશાસ્ત્રમાં આવરી લેવાયેલી સામગ્રીને
  ધ્યાનમાં લઈ શકે એવી વ્યવસ્થા કરવાની યોજના છે: પહેલેથી આવરી લેવાયેલી
  સામગ્રી કાઢી શકાય અને સાબિતીઓને વિધાનાત્મક ભાગનાં સંબંધિત સ્થાનોના
  સંદર્ભો આપીને ટૂંકાવી શકાય.
\end{editorial}

\olimport[introduction]{introduction}

\olimport[syntax-and-semantics]{syntax}

\olimport[syntax-and-semantics]{semantics}

\olimport[models-theories]{models-theories}

\olimport[proof-systems]{proof-systems}

\iftag{prfSC}{%
  \olimport[sequent-calculus]{sequent-calculus}
}{}

\iftag{prfND}{%
  \olimport[natural-deduction]{natural-deduction}
}{}

\iftag{prfTab}{%
  \olimport[tableaux]{tableaux}
}{}

\iftag{prfAX}{%
  \olimport[axiomatic-deduction]{axiomatic-deduction}
}{}

\olimport[completeness]{completeness}

\olimport[beyond]{beyond}

\OLEndPartHook

\end{document}
